\centerline{\bf \Huge Тренировочная олимпиада}
\centerline{\bf \Huge Муницип 7 класс}

\begin{flushright}
        \scriptsize{Каждая задача оценивается в 4 балла.}
\end{flushright}

\problem{7.1} Бегемотики Ася и Вася изначально весили одинаково. За год вес Аси увеличился в 3,5 раза, в следующий год - ещё в 2,5 раза, а в третий - ещё в 1,5 раза. При этом вес Васи в первый год увеличился в 4,5 раза, во второй - ещё в 3,5 раза, а в третий - ещё в 2,5 раза. Во сколько раз теперь Вася тяжелее Аси?

\problem{7.2} Петя написал на доске семизначное число. Для каждой пары соседних цифр этого числа Вася записал себе в тетрадку их сумму. Оказалось, что у Васи все числа различные. Какое максимальное число мог написать Петя?

\problem{7.3} На прямой в некотором порядке расположены пять различных точек $A, B, C, D$ и $E$. Известно, что $A B=4, B C=7, C D=11, D E=16$.

Чему равно наименьшее возможное расстояние между $A$ и $E$ ?

\problem{7.4} Малыш и Карлсон решили пробежать три круга по стадиону. Малыш бежал всю дистанцию с постоянной скоростью 4 км/ч.

Карлсон бежал каждый круг с постоянной скоростью. Первый круг он бежал с той же скоростью, что и Малыш. Затем он подкрепился вареньем и пробежал второй круг в шесть раз быстрее. После второго круга он тоже подкрепился вареньем, но понял, что объелся, и замедлился. С какой скоростью Карлсон бежал третий круг, если они с Малышом финишировали одновременно, но в течение забега Карлсон потратил на поедание варенья столько же времени, сколько и на бег? Ответ дайте в км/ч.

\problem{7.5} Маша, Даша и Саша загадали по числу от 1 до 9 , а затем сообщили эти числа друг другу. Оказалось, что все загаданные ребятами числа различны. После этого каждый из них произнёс по утверждению:
- Маша: «Сумма загаданных чисел делится на 4»
- Даша: «Если бы я могла загадывать числа больше 9, я бы загадала число в три раза больше, и тогда сумма загаданных увеличилась бы вдвое»
- Саша: «Все загаданные числа больше 2»

Напишите загаданные ими числа в любом порядке, если известно, что никто не из них не ошибся.